\documentclass[10pt,a4paper]{amsart}
\usepackage[margin=19mm]{geometry}
\usepackage{amsmath,amssymb,mathtools}
\usepackage[scr=rsfs,cal=euler]{mathalfa}
\usepackage{xfrac}
\usepackage[unicode,hidelinks,bookmarks=false]{hyperref}
\newcommand{\ie}{i.e.\ }
\newcommand{\cf}{cf.\ }
\newcommand{\resp}{respectively\ }
\newcommand{\Subsection}{\subsection}
\providecommand{\nobreakdash}{\nobreak\hbox{-}}
\setlength{\emergencystretch}{2em}
\multlinegap0pt
\usepackage[shortlabels]{enumitem}
\usepackage{amssymb}
\usepackage{mathtools}
\usepackage{tikz}
\usepackage{float}
\newtheorem{lemma}{Lemma}
\DeclareMathOperator{\Lip}{Lip}
\newcommand{\eps}{\varepsilon}
\title{\textbf{Logarithmic capacity and torsion in planar shape optimization}}
\author{David Ziener}
\date{Source: arXiv:2609.06205v1 (CC BY 4.0)}
\begin{document}
\maketitle

\noindent\emph{Source and licence.} \textbf{Logarithmic capacity and torsion in planar shape optimization}. Version arXiv:2609.06205v1 (CC BY 4.0), distributed by arXiv under the Creative Commons Attribution 4.0 International licence, http://creativecommons.org/licenses/by/4.0/, as recorded in the arXiv licence metadata for this exact version and shown on the abstract page https://arxiv.org/abs/2609.06205v1. Full licence notice: \texttt{licenses/arxiv-2609.06205-CC-BY-4.0.txt}.\par\smallskip

\noindent\textit{Editorial scope.} Counted unit: the article's lemma lem:MovingConvergence (source0.tex lines 832--852). The article's complete original proof of it is delivered with it (source0.tex lines 854--979, one \texttt{proof} environment of 126 lines). No mathematics is written, repaired or re-derived: the statement and the proof are the article's own lines, in the article's own notation, and no retained line is substituted or re-flowed.  Retained uncounted as the source-local context the proof needs: lines 802--812.  Nothing was omitted from any retained span, so the package carries no omission record.  No retained line contains a citation, so the wrapper carries no bibliography and no bibliography entry.  No retained line contains a figure environment, an image-inclusion command or drawing code, so no image is delivered and the package declares no asset dependency.  Editorial formatting only, in addition: an AMS \texttt{amsart} preamble that restates the article's own declarations and macros used by the retained lines, namely the environments \texttt{lemma} on separate counters, together with the standard packages those lines need.  The retained environments print in this document as Lemma 1 in order of appearance, whereas the article prints the same items with the numbering of its own full document.  The complete native source of the article as distributed in the arXiv e-print for this exact version is bundled unchanged as \texttt{sources/209504/source0.tex}.
    \begin{lemma}
        \label{lem:KernelExpansion}
        Suppose $h \in X$ with $\Lip(h)\leq \alpha <1$. Write $\lambda:=1-\alpha >0$. For each $q \in X$ with $\Lip(q)$ small enough one has the expansion
        \[
        K_{h+q}(x,y)-K_h(x,y)=B_{h,q}(x,y)+R_{h,q}(x,y),
        \]
        where
        \[
        |B_{h,q}|\leq C_\lambda \Lip(q),\quad |R_{h,q}|\leq C_\lambda\Lip(q)^2.
        \]
    \end{lemma}

    \begin{lemma}
    \label{lem:MovingConvergence}
    Let $g_n\to h$ and $\tilde h_n\to h$ in $X$, and assume that $g_n,\tilde h_n,h$ all lie eventually in a common uniformly bi-Lipschitz neighborhood. Let
    \[
        \mu_n:=\mu_{\tilde h_n},\qquad \mu:=\mu_h,
    \]
    and set
    \[
        \eta_n:=\mu_n\otimes\mu_n,\qquad \eta:=\mu\otimes\mu.
    \]
    Then
    \begin{equation*}
        \sup_{\|k\|_X\leq 1}
        \left|
            \int\int B_{g_n,k}\,d\mu_n d\mu_n
            -
            \int\int B_{g_n,k}\,d\mu d\mu
        \right|
        \to 0,\quad \text{as }n \to \infty.
    \end{equation*}
\end{lemma}

\begin{proof}
    We divide the proof into two steps. First we prove the convergence of measures, then we prove that this is enough to pass to the limit uniformly in the kernel integral.\\
    \textbf{Step 1: Convergence of measures.} We want to prove here that
    \begin{equation}
        \label{eq:ConvergenceOfMeas}
        \eta_n \rightharpoonup \eta.
    \end{equation}
    Since $\tilde h_n\to h$ in $X$ and all maps lie eventually in a common uniformly bi-Lipschitz neighborhood, integrating Lemma \ref{lem:KernelExpansion} gives
    \begin{equation}
        \label{eq:JBound}
        |J_{\tilde h_n}(\nu)-J_h(\nu)|
        \leq C\|\tilde h_n-h\|_X
    \end{equation}
    for the atomless probability measures used below. Hence, using the minimality of $\mu_n$ for $J_{\tilde h_n}$, we estimate
    \begin{align}
        \label{eq:MinSeq}
        J_h(\mu_n)
        &\leq J_{\tilde h_n}(\mu_n)+C\|\tilde h_n-h\|_X  \notag\\
        &\leq J_{\tilde h_n}(\mu)+C\|\tilde h_n-h\|_X \notag\\
        &\leq J_h(\mu)+2C\|\tilde h_n-h\|_X .
    \end{align}
    In particular, \eqref{eq:MinSeq} says that $(\mu_n)$ is a minimizing sequence for $J_h$. Up to taking subsequences, $\mu_n\rightharpoonup \tilde\mu$. By lower semicontinuity of the logarithmic kernel, $\tilde\mu$ is minimizing for $J_h$. By uniqueness it follows that $\tilde\mu=\mu$. Since the convergence holds for all subsequences, we have shown that $\mu_n\rightharpoonup\mu$. The convergence \eqref{eq:ConvergenceOfMeas} follows by compactness of $E$.\\
    \textbf{Step 2: Uniform convergence of kernels.} We want to prove
    \begin{equation}
        \label{eq:KernelConv}
        \sup_{\|k\|_X\leq 1}
        \left|
            \int\int B_{g_n,k}\,d\mu_n d\mu_n
            -
            \int\int B_{g_n,k}\,d\mu d\mu
        \right|
        \to 0.
    \end{equation}
    Since all the maps lie in a common uniformly bi-Lipschitz neighborhood, there exists $\lambda>0$ such that
    \[
        |\Phi_{g_n}(x)-\Phi_{g_n}(y)|\geq \lambda |x-y|
    \]
    for all $n$ sufficiently large and all $x,y\in E$. Hence as in the proof of Lemma \ref{lem:KernelExpansion}
    \begin{equation}
        \label{eq:KernelUniformBound}
        |B_{g_n,k}(x,y)|
        \leq C_\lambda,
        \qquad x\neq y,\quad \|k\|_X\leq 1.
    \end{equation}
    Fix $\eps>0$. Since $\eta$ gives no mass to the diagonal, we can choose $\delta>0$ such that
    \begin{equation*}
        \eta(F_{2\delta})<\eps,
        \qquad
        F_\rho:=\{(x,y)\in E\times E: |x-y|\leq \rho\}.
    \end{equation*}
    From \eqref{eq:ConvergenceOfMeas} and because $F_{2\delta}$ is closed, we have
    \begin{equation}
        \label{eq:LimSupSmall}
        \limsup_{n\to\infty}\eta_n(F_{2\delta})
        \leq \eta(F_{2\delta})
        <\eps.
    \end{equation}
    Let $\chi_\delta$ be a smooth cutoff such that $\chi_\delta\equiv 0$ on $F_\delta$ and $\chi_\delta\equiv 1$ on $E\times E\setminus F_{2\delta}$. We claim that the family
    \[
        \mathcal F_\delta
        :=
        \{\chi_\delta B_{g_n,k}: n\in\mathbb N,\ \|k\|_X\leq 1\}
    \]
    is equicontinuous and uniformly bounded on $E\times E$. Uniform boundedness follows immediately from \eqref{eq:KernelUniformBound}. For equicontinuity close to the diagonal, the claim follows from the cutoff. In the set $E\times E\setminus F_\delta$, note that
    \[
        (x,y,a,b)\mapsto
        -\frac{(\Phi_{g_n}(x)-\Phi_{g_n}(y))\cdot(a-b)}
        {|\Phi_{g_n}(x)-\Phi_{g_n}(y)|^2}
    \]
    is uniformly Lipschitz, because $g_n\to h$ in $X$ and we stay a positive distance away from the diagonal. Moreover, the maps $k$ with $\|k\|_X\leq 1$ are uniformly bounded and uniformly Lipschitz. This yields the claim.\\
    Hence, from Arzelà-Ascoli, the family $\mathcal F_\delta$ is totally bounded in $C(E\times E)$. Thus, we can pick a finite set $\{\psi_1,\ldots,\psi_N\}\subset C(E\times E)$ such that for every $n$ and every $k$ with $\|k\|_X\leq 1$, there is an index $j(n,k)\in\{1,\ldots,N\}$ satisfying
    \[
        \|\chi_\delta B_{g_n,k}-\psi_{j(n,k)}\|_{L^\infty}<\eps.
    \]
    From the convergence \eqref{eq:ConvergenceOfMeas}, we have
    \[
        \max_{1\leq j\leq N}
        \left|
            \int \psi_j\,d\eta_n
            -
            \int \psi_j\,d\eta
        \right|
        \to 0.
    \]
    Hence,
    \[
        \sup_{\|k\|_X\leq 1}
        \left|
            \int \chi_\delta B_{g_n,k}\,d\eta_n
            -
            \int \chi_\delta B_{g_n,k}\,d\eta
        \right|
        \leq 2\eps+o(1).
    \]
    Finally, using the properties of the cutoff, the uniform bound \eqref{eq:KernelUniformBound}, and \eqref{eq:LimSupSmall}, we get
    \begin{align*}
        &\sup_{\|k\|_X\leq 1}
        \left|
            \int B_{g_n,k}\,d\eta_n
            -
            \int B_{g_n,k}\,d\eta
        \right| \\
        &\leq
        \sup_{\|k\|_X\leq 1}
        \left|
            \int \chi_\delta B_{g_n,k}\,d\eta_n
            -
            \int \chi_\delta B_{g_n,k}\,d\eta
        \right| \\
        &\qquad
        +
        \sup_{\|k\|_X\leq 1}
        \left|
            \int (1-\chi_\delta) B_{g_n,k}\,d\eta_n
            -
            \int (1-\chi_\delta) B_{g_n,k}\,d\eta
        \right| \\
        &\leq
        2\eps+o(1)
        +
        C_\lambda\bigl(\eta_n(F_{2\delta})+\eta(F_{2\delta})\bigr) \\
        &\leq
        2\eps+C_\lambda\eps+o(1).
    \end{align*}
    Since $\eps>0$ was arbitrary, this completes the proof of \eqref{eq:KernelConv} and hence the proof of the lemma.
\end{proof}

\end{document}
